%=============================================================================
\section{Theoretical background}
\label{sec:theory}
%=============================================================================

\subsection{Anticipatory governance and the demand-side Collingridge dilemma}

Anticipatory governance designates the capacity of a society to act on
emerging sociotechnical change while that change can still be shaped
\citep{barben2008,guston2014}. It rests on three coupled capacities: foresight,
engagement with affected publics, and integration of both into decision
routines. \citet{fuerth2009} adds a systems requirement: foresight governs only
when it is wired into the policy process and fed back as evidence accumulates.
Applied to planning, the same logic replaces prediction with the preparation of
responses across a range of plausible futures \citep{quay2010}.

The underlying problem is the dilemma of control \citep{collingridge1980}.
Early in the diffusion of a technology, intervention is cheap but its
consequences cannot be known; once consequences are known, the technology is
entrenched and intervention is costly. The dilemma is usually stated for the
supply side, that is, for the impacts of an artefact. Regulatory mandates
create a demand-side variant. Regulation~2025/40 fixes the technology and the
timeline, so the unknown is no longer what the technology does but how
heterogeneous publics will respond to it. Before the 2030 targets bind,
pricing, labelling, and communication can still be adjusted at low cost, yet
consumer response has not been measured. Once response becomes observable in
market data, firms have committed capital and any loss of trust has already
occurred. The policy-before-data problem is this demand-side Collingridge
dilemma, and it defines the window in which anticipatory instruments have
value.

\subsection{Risk mitigation under regulatory transition in the recycling sector}

Governance under such conditions is necessarily tentative: provisional,
revisable, and designed to learn \citep{kuhlmann2019}. Adaptive planning
operationalises this through commitments staged against monitored signposts
\citep{walker2013}, and transition studies show that a mandate performs as part
of a policy mix, in which flanking instruments determine whether the core
instrument retains legitimacy \citep{rogge2016}. Risk mitigation in a
regulatory transition therefore consists of three decisions: where to intervene
first, with which flanking instrument, and when to revise.

The appropriate strategy depends on the kind of risk. \citet{klinke2002}
distinguish risk-based, precaution-based, and discourse-based management, the
last being required when the problem is ambiguity rather than lack of data.
Three sources of indeterminacy must accordingly be kept apart in the present
case: \emph{variability} (consumers and segments differ), \emph{epistemic
uncertainty} (the outcomes have not been measured in the target population and
must be inferred from other populations), and \emph{ambiguity} (actors disagree
on what counts as safe, fair, or acceptable). A quantitative instrument can
describe the first, can address the second only under explicit transfer
assumptions, and cannot resolve the third.

Recycled-content food-contact packaging sharpens each decision. The risk object
is double. Technical safety is governed \emph{ex ante}: recycling processes
require authorisation under Regulation (EU)~2022/1616 following assessment by
the European Food Safety Authority \citep{eu2022}. Perceived safety is not
governed at all. Consumers associate recycled material with contamination; the
association lowers acceptance independently of certified safety and is
strongest for products in contact with the body or with food
\citep{magnier2019,meng2021}. The binding targets remove the option of waiting,
and perception varies across demographic and attitudinal segments
\citep{ruokamo2022,zwicker2020}. Mitigation is consequently a targeting
problem: the flanking measures that matter (traceability, certification marks,
safety-framed communication, staged rollout) can be effective only if directed
at the segments where the perception gap is largest.

\subsection{Design requirements for an anticipatory instrument}
\label{sec:requirements}

Four requirements follow. \emph{(R1) Pre-observational operation}: the
instrument must produce estimates for variables not yet measured in the target
population. Measurement does not close the Collingridge window; the window
narrows as commitments, infrastructure, and practices consolidate, and an
instrument that waits for measurement acts when adjustment has become costly.
\emph{(R2) Segment-level resolution}: because mitigation is a targeting
problem, outputs must be distributions conditional on demographic and
attitudinal profiles, with explicit uncertainty. \emph{(R3) Updatability}:
tentative governance requires that new evidence revise the estimates without
rebuilding the instrument. \emph{(R4) Accessibility}: the engagement and
integration capacities require that actors without statistical training,
including SMEs and regulators, can query the instrument and inspect its
assumptions.

Delphi panels and scenario methods satisfy R1, and iterative designs revise
their outputs between rounds, but they do not usually provide quantitative
conditional resolution and formal probabilistic updating together. Statistical
forecasting and synthetic-data generators trained on a single dataset satisfy
R2 only for variables already observed, and so fail R1.
Sections~\ref{sec:system} and~\ref{sec:framework} present BVS as an instrument
built to satisfy all four.

Within foresight methodology, BVS belongs to what \citet{bossler2026} call
\emph{synthetic foresight}: a mode in which the participants, populations, and
evidence bases of a foresight exercise are computationally generated rather
than recruited. They position it as the sixth column of the typology of
approaches to the future of \citet{bas2004} and \citet{bas2015}
(Table~\ref{tab:typology}), with synthetic Delphi panels, language-model
virtual surveys, and Bayesian virtual surveys as its methods. Their evaluation
of a 60-agent synthetic Delphi panel over a knowledge graph of 25{,}823
circular-economy publications bears directly on R1--R4: the panel converged
under every one of eight evidence configurations, with reductions of 81--91\%
in the coefficient of variation, yet convergence was independent of coverage of
the field's research communities, and coverage and ranking accuracy against
observed 2021--2024 publication growth depended on the type of grounding.
Agreement among synthetic respondents is therefore not a validity criterion.
The same holds for BVS: a coherent synthetic population says nothing about
whether the transferred evidence is adequate, so validity has to be built into
the evidence architecture (Section~\ref{sec:data}) and tested against an
external record (Section~\ref{sec:step5}).

\begin{table}[htbp]
\centering
\caption{Approaches to the future: a typology. Columns 1--5 build on Bas [1]
and Bas and Guill\'o [2], drawing on Habermas's four knowledge interests. The
integrated typology and its Synthetic extension were conceptualized and
developed by Bas and Bossler. The sequence should be read as a series of
substrate transitions rather than as discrete schools of thought.
Reproduced without modification from \citet[Table~1]{bossler2026} under a
CC~BY-NC-ND~4.0 licence; [1] and [2] refer to \citet{bas2004} and
\citet{bas2015}.}
\label{tab:typology}
\scriptsize
\setlength{\tabcolsep}{4pt}
\renewcommand{\arraystretch}{1.15}
\begin{tabularx}{\textwidth}{>{\itshape}l*{6}{>{\raggedright\arraybackslash}X}}
\toprule
\normalfont & \textbf{Supernatural} & \textbf{Hermeneutic} & \textbf{Technical} &
  \textbf{Emancipatory} & \textbf{Participatory} & \textbf{Synthetic} \\
\normalfont & & & & & & \textbf{\emph{(proposed)}} \\
\midrule
Origin & Primitive societies & Secular society & Industrial society &
  Post-industrial & Liquid society & Digital society \\
Paradigm shift & Old regime & Nation-state & Mass capitalism & Glocalization &
  Social networks & Foundation models \\
Key point & Religion & Ideology & Product & Knowledge & People/\allowbreak Experience & Data \\
Dominant subsystem & Culture & Politics & Economy & Technology & Ecosystem &
  Virtual \\
Foundation & Emotional & Emotional & Rational & Rational & Rational + Emotional &
  Emotional + Rational \\
Method & God's orders & Utopias & Trend forecasting & Foresight/\allowbreak Delphi &
  Creative futures & Augmented intelligence \\
Attitude & Alienation & Alienation & Adaptation & Anticipation & Innovation &
  Augmentation \\
\bottomrule
\end{tabularx}
\end{table}

Two boundary conditions qualify this claim. First, an instrument that meets
R1--R4 contributes to the \emph{foresight} capacity only. Synthetic respondents
are computed from evidence about other people; they do not participate,
consent, or deliberate, and they cannot stand in for engagement. Second, an
instrument used in rule-making is itself a technology to be governed.
Responsible-innovation criteria of anticipation, reflexivity, inclusion, and
responsiveness \citep{stilgoe2013} apply to BVS as much as to the packaging
transition: who selects the sources, who validates the mapping between survey
questions, who can contest an estimate, and who answers for its use.
Section~\ref{sec:implications} returns to these conditions.

